\section{What the Soul Is Doing}

Ask a catechism class what happens when a Christian dies and the answer will come back promptly: the soul goes to heaven. Ask where the Church defined that and the room goes quiet, and rightly, because the answer is more interesting than the question expects. Something very like it \emph{was} defined, in 1336, by a pope whose predecessor had preached the opposite from his own pulpit and retracted it on his deathbed.

Getting to that requires first admitting how little the older evidence says.

\subsection{The reticence of the Hebrew Scriptures}

The Old Testament does not, for the most part, hold out a conscious and blessed life immediately after death. It says the opposite often enough that the Fathers had to work to accommodate it.

The psalmist bargains with God on the assumption that the dead are useless to him: ``For there is no one in death, that is mindful of thee: and who shall confess to thee in hell?''\endnote{Psalm 6:6 (Vulgate numbering), Douay--Rheims, read 2026-07-26 in the tracked Challoner verse text. The Latin \latin{in inferno} renders the Hebrew \latin{sheol}; the Douay's ``hell'' here is not the hell of the damned but the common abode of the dead, and reading it otherwise produces nonsense. Augustine cites this verse at \emph{De civitate Dei} XIII.11 for a different purpose---to show that Scripture speaks of being \emph{in} death---and the essay records that its own use of the verse is not his.} Job pushes further, into flat denial: ``But man when he shall be dead, and stripped and consumed, I pray you where is he? \ldots{} So man when he is fallen asleep shall not rise again; till the heavens be broken, he shall not awake, nor rise up out of his sleep.'' And then, within four verses, the question that the whole later tradition is an answer to: ``Shall man that is dead, thinkest thou, live again? all the days in which I am now in warfare, I expect until my change come.''\endnote{Job 14:10, 12, 14, Douay--Rheims, read as above. The chapter opens with the famous ``Man born of a woman, living for a short time, is filled with many miseries'' (14:1). Verse 13---``Who will grant me this, that thou mayst protect me in hell, and hide me till thy wrath pass, and appoint me a time when thou wilt remember me?''---is a request rather than an assertion, and the essay does not read it as a doctrine of the intermediate state.} Ecclesiastes ends its long meditation on decay with a sentence that says something and refuses to say more: ``And the dust return into its earth, from whence it was, and the spirit return to God, who gave it.''\endnote{Ecclesiastes 12:7, Douay--Rheims, read as above; the preceding verses (12:1--6) are the well-known figure of the failing body. The verse is quoted by CCC 1007 to make exactly the point made here---that death lends urgency to a measured life---and the Catechism draws no further inference from it.} The spirit returns to God; and there the sentence stops.

The change comes late and in Greek. ``But the souls of the just are in the hand of God, and the torment of death shall not touch them. In the sight of the unwise they seemed to die: and their departure was taken for misery: and their going away from us, for utter destruction: but they are in peace.''\endnote{Wisdom 3:1--3, Douay--Rheims, read as above. The passage is explicitly correcting an inference---what \emph{seemed} to be the case to ``the unwise''---which suggests that the older and bleaker reading was the ordinary one among its first readers. Wisdom is a first-century-\textsc{bc} Alexandrian book preserved in Greek and received in the Catholic canon; the essay does not enter the canonical dispute.}

The Catechism does not disguise this trajectory. ``God revealed the resurrection of the dead to his people progressively,'' it says, and locates the decisive moment in the Maccabean martyrs.\endnote{CCC 992, Holy See English web text, read 2026-07-26; the paragraph quotes 2 Maccabees 7:9 and 14. CCC 993 adds that ``the Pharisees and many of the Lord's contemporaries hoped for the resurrection,'' and that the Sadducees denied it. That a doctrine was revealed progressively is a claim about the mode of revelation, not a concession that the earlier books are erroneous; the essay records the development as the Catechism states it.} That is an honest description of the record, and worth setting down plainly, because a reader who has been told that the Bible teaches an immortal soul going to its reward will be startled the first time he reads Job 14 and will be right to be. The doctrine is genuinely there in the end. It is not there from the beginning, and the tradition has never pretended otherwise.

\subsection{The New Testament passages, each in its own place}

The relevant New Testament texts are few, and each has been made to carry more than it says by being read through the others. Read singly they are more modest and more useful.

To the thief: ``Amen I say to thee: This day thou shalt be with me in paradise.''\endnote{Luke 23:43, Douay--Rheims, read 2026-07-26 in the tracked Challoner verse text. The reader should note that the punctuation---whether ``Amen I say to thee this day'' or ``This day thou shalt be with me''---has been contested, since the ancient manuscripts are unpunctuated; the Douay follows the Latin tradition's placement, which the Catholic reading assumes. The essay flags the question and does not adjudicate it, having examined neither the Greek witnesses nor the punctuation history.} The word doing the work is \emph{this day}. Whatever paradise is, the promise attaches an immediate term to it, and the tradition has always read the verse that way.

To the Corinthians, Paul writes the most-quoted and least-read passage on the subject. Its centre is not the confidence at the end but the recoil in the middle: ``For we also, who are in this tabernacle, do groan, being burthened; because we would not be unclothed, but clothed upon, that that which is mortal may be swallowed up by life.''\endnote{2 Corinthians 5:4, Douay--Rheims, read as above; the surrounding verses are 5:1 (``if our earthly house of this habitation be dissolved, that we have a building of God, a house not made with hands, eternal in heaven''), 5:6 (``while we are in the body we are absent from the Lord'') and 5:8 (``we are confident and have a good will to be absent rather from the body and to be present with the Lord''). The reading offered here---that verse 4 states a preference \emph{against} the unclothed condition---follows the plain sense of ``we would not be unclothed, but clothed upon'' and is offered as such; commentators differ about the referent of the clothing imagery, and the essay does not survey them.} Paul does not want to be \emph{unclothed}. He wants to be \emph{further clothed}. Given the choice he would skip the disembodied condition altogether and go straight from mortal to swallowed-up-by-life. Only after registering that preference does he say the sentence everyone quotes---that he would rather be absent from the body and present with the Lord---and the order of the two is the whole point. Presence with the Lord is worth being unclothed for. It does not make being unclothed desirable.

To John on Patmos: ``I saw under the altar the souls of them that were slain for the word of God\ldots{} And white robes were given to every one of them one; And it was said to them that they should rest for a little time till their fellow servants and their brethren\ldots{} should be filled up.''\endnote{Apocalypse 6:9--11, Douay--Rheims, read as above. The essay takes the vision as an apocalyptic image and not as a topography; the point drawn from it---that the souls are conscious, are answered, and are told to wait---is on the surface of the text and does not require decoding the imagery.} Conscious; audible; given something; and told to wait. It is a small datum and an exact one: whatever the state is, it has a duration and it is not the end.

And the parable of the rich man and Lazarus, which is where most Catholic imagination about the interval actually comes from, and which is a parable. Its furniture---the bosom of Abraham, the fixed chasm, the sight across the gap, the thirst---has been enormously influential and was not offered as a survey.\endnote{Luke 16:22--26, Douay--Rheims, read as above. CCC 1021 cites it (with Luke 23:43) as among the New Testament texts that ``speak of a final destiny of the soul,'' which is a careful formulation: the Catechism draws from the parable the fact of a differentiated destiny, not the details of the scene. The caution in the text above is this essay's and is directed at the popular use of the parable's imagery as description.}

That is the evidence. It establishes that something of the person persists, that its condition is differentiated according to how the person lived, that it is conscious, and that it is waiting. It does not establish what the waiting is like, and the New Testament nowhere describes it.

\subsection{The philosophers' problem: a part that is not a person}

What the scholastics added was not more information but an analysis of what kind of thing could be in such a state, and the analysis is uncomfortable in a way that has largely been forgotten.

Aquinas asks whether the separated soul can understand anything at all, and the difficulty he sets himself is severe: on his own account the soul understands by turning to images supplied through the body, and death removes the body. His answer is that the soul's mode of operating follows its mode of being, and that both change: separated, it understands by turning to intelligible objects directly, as the angels do. And then the sentence that costs him something: ``it is as natural for the soul to understand by turning to the phantasms as it is for it to be joined to the body; but to be separated from the body is not in accordance with its nature, and likewise to understand without turning to the phantasms is not natural to it.''\endnote{ST I, q.~89, a.~1, corpus; English Dominican translation, read 2026-07-26 at New Advent. The corpus argues expressly against the Platonic escape---``Did this not proceed from the soul's very nature, but accidentally through its being bound up with the body, as the Platonists said, the difficulty would vanish''---and rejects it on the ground that in that case ``the union of soul and body would not be for the soul's good.'' Aquinas's comparison is a light body outside its proper place. The corpus goes on to argue that the separated soul's knowledge, though of a nobler kind, is in it ``imperfect'' and ``of a general and confused nature'' relative to what higher substances do with the same mode.}

The separated soul, on the most influential account the Catholic tradition possesses, is not liberated. It is displaced. It is operating in a mode above its station and below its capacity, like a hand playing an instrument it was not shaped for.

And it is not the person. Commenting on 1 Corinthians 15, Aquinas argues that a denial of the bodily resurrection cannot be repaired by an appeal to the immortality of the soul, and he gives two reasons. The first is that a soul permanently without a body would make the unnatural condition permanent and the natural one transient, which is metaphysically backwards. The second is the one that has echoed: ``since the soul is a part of the body of a man, it is not the whole man, and my soul is not I; so that although the soul obtains salvation in another life, nevertheless not I or any man.''\endnote{Thomas Aquinas, \emph{Super I Epistolam B. Pauli ad Corinthios lectura}, cap.~15, lect.~2: ``\latin{anima autem cum sit pars corporis hominis, non est totus homo, et anima mea non est ego; unde licet anima consequatur salutem in alia vita, non tamen ego vel quilibet homo.}'' Latin read 2026-07-26 in the Corpus Thomisticum; the English above is a working gloss. A textual caution is required and is usually omitted when this sentence is quoted: the commentary on 1 Corinthians survives in several states, and chapters 11--16---which contain this passage---are transmitted in the \latin{reportatio vulgata}, a reported text, not in a portion written in Thomas's own hand. The Corpus Thomisticum labels the section accordingly. The sentence is therefore excellent evidence of the Thomistic school's teaching and weaker evidence of Thomas's own words than its fame suggests. The preceding argument in the same passage---``the soul is naturally united to the body, and is separated from it against its nature and \latin{per accidens}; whence the soul stripped of the body, so long as it is without the body, is imperfect''---is on the same footing.}

Set that beside the definition of death from which this essay started, and the intermediate state acquires its proper strangeness. Death separates soul and body. The soul survives. The soul is not the man. Therefore what survives death is not, strictly, the person who died---and yet the Church prays to the saints by name, and offers Masses for the dead by name, and the whole practice presupposes that somebody is there to be addressed.

\subsection{What the Church actually defined, and the crisis that made it define it}

The doctrine that the blessed see God immediately, before the resurrection, is not a piece of folk piety that hardened. It is a papal definition, and it was provoked by a pope.

In the winter of 1331--1332 John XXII preached at Avignon a series of sermons proposing that the souls of the just, before the general resurrection, do not yet see the divine essence but rest under the altar in the humanity of Christ, receiving the vision only after they take their bodies again. The proposal caused an uproar. On 3 December 1334, the day before he died, he issued \latin{Ne super his}, declaring his mind: ``we confess and believe that the purified souls separated from their bodies are in heaven, in the kingdom of heaven and paradise, and are gathered with Christ in the company of the angels, and see God and the divine essence face to face clearly, according to the common law, \emph{in so far as the state and condition of a separated soul allows}.''\endnote{John XXII, bull \latin{Ne super his} (3 December 1334), DS 990--991: ``\latin{Fatemur siquidem et credimus, quod animae purgatae separatae a corporibus sunt in caelo, caelorum regno et paradiso et cum Christo in consortio angelorum congregatae et vident Deum de communi lege ac divinam essentiam facie ad faciem clare, in quantum status et condicio compatitur animae separatae.}'' Latin read 2026-07-26 at the patristica.net presentation of the Latin Enchiridion; the English is a working gloss and the emphasis is this essay's. The document's own opening explains its purpose: that what had been said ``both by Us and by several others in Our presence'' about whether purified souls see the divine essence before resuming their bodies should not reach the ears of the faithful otherwise than as he intended it. CCC 1022 cites this act in its footnotes alongside \emph{Benedictus Deus}. The essay reports the sermons as the occasion of the controversy on the basis of the bull's own description of the disputed question, and does not claim to have examined the sermons.}

Fourteen months later his successor closed the question. Benedict XII's constitution \latin{Benedictus Deus} of 29 January 1336 defines that the souls of the saints and of the faithful departed in whom nothing remained to be purged---and of those purged after death, and of baptized children dying before the use of reason---``immediately after their death and the aforesaid purification\ldots{} \emph{even before the resumption of their bodies and the general judgment}\ldots{} have been, are, and will be in heaven, in the kingdom of heaven and the heavenly paradise with Christ''; and that they ``have seen and do see the divine essence with an intuitive and even facial vision, with no creature intervening in the character of an object seen, but the divine essence showing itself to them immediately, nakedly, clearly, and openly.''\endnote{Benedict XII, constitution \latin{Benedictus Deus} (29 January 1336), DS 1000: ``\latin{Hac in perpetuum valitura Constitutione auctoritate Apostolica diffinimus: quod\ldots{} mox post mortem suam et purgationem praefatam\ldots{} etiam ante resumptionem suorum corporum et iudicium generale\ldots{} fuerunt, sunt et erunt in caelo\ldots{} viderunt et vident divinam essentiam visione intuitiva et etiam faciali, nulla mediante creatura in ratione obiecti visi se habente, sed divina essentia immediate se nude, clare et aperte eis ostendente.}'' Latin read 2026-07-26 at the patristica.net presentation of the Latin Enchiridion; the emphasis is this essay's. The central clause is quoted in the official English of CCC 1023: ``these souls have seen and do see the divine essence with an intuitive vision, and even face to face, without the mediation of any creature.'' The register is definitive: \latin{diffinimus}, in a constitution declared to be of perpetual force.} And the other side of it, in the same act: ``the souls of those who depart in actual mortal sin descend immediately after their death into hell, where they are tormented by infernal punishments.''\endnote{\emph{Benedictus Deus}, DS 1002: ``\latin{Diffinimus insuper, quod secundum Dei ordinationem communem animae decedentium in actuali peccato mortali mox post mortem suam ad inferna descendunt, ubi poenis infernalibus cruciantur, et quod nihilominus in die iudicii omnes homines `ante tribunal Christi' cum suis corporibus comparebunt, reddituri de factis propriis rationem.}'' Latin read as above; the closing quotation is 2 Corinthians 5:10. The third volume of this series treats the punishment; the point taken here is only the timing.}

Two features of that definition deserve more attention than they usually get.

First, John XXII's careful qualifier---\emph{in so far as the state and condition of a separated soul allows}---does not appear in Benedict's text. Where the earlier document hedges the vision by the limits of the seer, the later one insists on the unmediated character of what is seen and says nothing about the seer's incapacity. The difference is not nothing, and a reader who wants to know what has been defined should notice that the defined form is the stronger one.

Second, and more important for this essay: what \emph{Benedictus Deus} defines is a \emph{when} and a \emph{what}, and nothing else. It fixes the timing---immediately, before the bodies are resumed, before the general judgment. It fixes the object and its immediacy---the divine essence, no creature mediating. It says nothing whatever about the mode of the separated soul's knowledge, nothing about its relation to time, nothing about whether it perceives events on earth, nothing about the texture of the interval. Everything this essay reported from Aquinas about displacement and incompleteness stands untouched by the definition, and everything popular piety has added about the dead watching over us is likewise untouched---which is to say, unsupported by it.

The Second Vatican Council states the resulting picture without adding to it: some of Christ's disciples ``are exiles on earth, some having died are purified, and others are in glory beholding `clearly God Himself triune and one, as He is'\,''---three conditions in one communion, the middle one temporary.\endnote{\emph{Lumen gentium} 49, Holy See English web text, read 2026-07-26; the internal quotation is from \emph{Benedictus Deus}. LG 50 grounds the Church's prayer for the dead in the same communion, citing 2 Maccabees 12:46 (``because it is a holy and wholesome thought to pray for the dead that they may be loosed from their sins''). The second and third volumes of this series treat judgment and its outcomes; the purification is not developed here.}

\subsection{The modern pressure, and where the Church drew the line}

In the second half of the twentieth century a serious theological proposal put the whole interval in question: that the resurrection occurs for each person at the moment of his own death, so that there is no waiting and no separated soul to wait---the ``resurrection in death.'' The motivation is not frivolous, and it is worth stating at its strongest, because its best argument is a Thomistic one. If my soul is not I, and if a soul without a body is in a condition contrary to its nature, then an interval of indefinite length in which the unnatural condition is the actual one looks like a defect in the arrangement rather than a feature of it. Resurrection in death removes the defect at a stroke.

The Congregation for the Doctrine of the Faith addressed the resulting confusion in 1979, in a letter to the world's bishops. Its method is worth noting: it does not name a theologian, does not pronounce a censure, and does not resolve the speculative question. It restates what may not be given up. ``The Church affirms that a spiritual element survives and subsists after death, an element endowed with consciousness and will, so that the `human self' subsists. To designate this element, the Church uses the word `soul', the accepted term in the usage of Scripture and Tradition.'' And: ``the Church looks for `the glorious manifestation of our Lord, Jesus Christ', believing it to be distinct and deferred with respect to the situation of people immediately after death.''\endnote{Congregation for the Doctrine of the Faith, \emph{Recentiores episcoporum Synodi}, letter on certain questions concerning eschatology (17 May 1979), points 3 and 5; Holy See English web text, read 2026-07-26. Point 4 is also to the purpose: ``The Church excludes every way of thinking or speaking that would render meaningless or unintelligible her prayers, her funeral rites and the religious acts offered for the dead. All these are, in their substance, \latin{loci theologici}.'' The letter was approved by John Paul II and ordered published; it is authoritative non-definitive teaching of a curial congregation, not a definition, and its own register is that of recalling what the Church already holds rather than of deciding a disputed question.}

That is a boundary rather than a verdict, and the difference matters. Point 4 shows the reasoning: the Church's practice---requiem, suffrage, funeral---is itself a theological source, and a theory which made those acts unintelligible would be refuted by them before any argument was needed. One does not pray for someone who has already been raised and judged and glorified. The letter's argument is that the Church's knees have a vote.

So the state of the question, stated by weight, is this. That the purified see God immediately and before the resurrection, and that the damned descend immediately: definitive, \emph{Benedictus Deus}. That a conscious subject---``the human self''---survives, and that the general resurrection is deferred with respect to it: authoritative teaching, restated in 1979 against a live pressure. That the separated soul's mode of knowing is altered and its condition is contrary to its nature: common theology of the Thomistic school, not defined. That my soul is not I: the same, with a textual caveat about the report in which the sentence survives. That the interval has a particular texture, that the dead observe the living, that time passes there as here: not taught at any level, and the popular confidence about it has no support in any document quoted in this essay.

Whatever else the interval is, it begins with something the Church has been more precise about than about anything else in this territory: the moment of accounting.
